\documentclass{article}

\usepackage{amsmath,amssymb}
\usepackage{xurl}
\usepackage[hidelinks]{hyperref}
\setlength{\emergencystretch}{2em}

\input{./command}

\title{The Nonzero-Coefficient Convention for CPSV16 Canonical Forms}
\author{}
\date{2026-08-23}

\begin{document}
\maketitle
\gapnote{local-correction}{resolved}

\section{The convention}

Cirac--P\'erez-Garc\'ia--Schuch--Verstraete construct the canonical form by
listing the invariant blocks of the completely positive map and normalizing
each block separately \cite[lines 214--225]{Cirac2016MPDO_arXiv}. The
construction therefore produces precisely the blocks that occur in the
tensor. Equation~\texttt{eq:II\_CF1} of
Ref.~\cite[lines 238--244]{Cirac2016MPDO_arXiv} defines
\begin{align}
    A^i
        &= \bigoplus_{k=1}^r \mu_k A_k^i.
    \label{eq:cpsvbnt_canonical_form_decomposition}
\end{align}
The displayed definition does not repeat that every $\mu_k$ is nonzero.
However, the subsequent normalization $|\mu_k|\leq 1$, with at least one
coefficient equal to one, presupposes this condition
\cite[line 246]{Cirac2016MPDO_arXiv}: a summand with coefficient zero
contributes nothing to any positive-length matrix product vector and is not a
block of the decomposition.

The formalization reads the canonical form with every coefficient $\mu_k$
nonzero. This requirement is part of the canonical-form data.\footnote{Formal
field: \leanid{MPSTensor.CPSVCanonicalFormData.weights_ne_zero} in
\doclink{TNLean/MPS/CanonicalForm/Definitions.lean}.} A basis of normal tensors
in the sense of Definition~2.4 of Ref.~\cite{Cirac2016MPDO_arXiv} is a family
of normal tensors whose matrix product vectors span those of $A$ at every
positive length and are eventually linearly independent. After
gauge-phase-equivalent blocks are identified, the blocks of a canonical form
constitute such a basis, and every basis element occurs with nonzero
coefficient.

\section{Consequences}

Under this convention, the following source statements are formalized as
stated, with the indicated blueprint nodes.

\begin{itemize}
    \item Proposition~2.7 of
          Ref.~\cite[lines 278--280 and 1135--1146]
          {Cirac2016MPDO_arXiv} states that a family is a basis of normal
          tensors for a canonical form if and only if every representative is
          normal, every block is gauge-phase equivalent to a representative,
          and distinct representatives are not gauge-phase equivalent. The
          corresponding blueprint node is
          \path{thm:cpsv_bnt_characterization}.\footnote{Formal declaration:
          \leanid{MPSTensor.isCPSVBasisOfNormalTensors_iff_canonicalForm_covered_and_minimal}.}

    \item The uniqueness statement in
          Ref.~\cite[line 1148]{Cirac2016MPDO_arXiv} says that two bases of
          normal tensors for the same positive-length matrix product vector
          family admit a bijective gauge-phase matching. It follows by
          applying Proposition~2.7 of
          Ref.~\cite{Cirac2016MPDO_arXiv} in both directions. The statement is
          formalized for basis-of-normal-tensors sector presentations, in
          which every basis element carries a nonzero coefficient. The
          corresponding blueprint node is
          \path{thm:cpsv_bnt_presentation_uniqueness}.\footnote{Formal
          declaration:
          \leanid{MPSTensor.IsBNTSectorPresentation.equiv_of_sameMPV₂Pos}.}

          The literal Definition~2.4 family predicate\footnote{Formal
          declaration: \leanid{MPSTensor.IsCPSVBasisOfNormalTensors}.} remains
          available as an interface and carries no uniqueness claim. For a
          one-block tensor $A_0$ with coefficient $1$, both $(A_0)$ and
          $(A_0,A_1)$ satisfy that family predicate if the coefficient of
          $A_1$ is set to zero at every length. The second family contains a
          zero-coefficient member. The convention excludes that member from
          every basis presentation, so this construction is not a
          counterexample to the formalized uniqueness statement.

    \item The proportional and equal fundamental theorems and their unitary
          refinements are Theorem~2.10, Corollary~2.11, and Corollary~A.6 of
          Ref.~\cite[lines 349--361 and 1167--1199]
          {Cirac2016MPDO_arXiv}. Their blueprint nodes are
          \path{thm:cpgsv_multiblock_ft_source},
          \path{thm:cpgsv_equal_case_source},
          \path{cor:sector_bnt_proportional_unitary_sector_match}, and
          \path{cor:sector_bnt_equal_unitary_global_gauge}.

    \item Lemma~A.5 of
          Ref.~\cite[lines 1155--1163]{Cirac2016MPDO_arXiv} states that
          finite-range power sums of nonzero lists determine equal lengths and
          equal multisets. The lemma applies to the nonzero coefficients of a
          canonical form. Its blueprint node is
          \path{thm:bounded_power_sum_multiset}.

    \item Corollary~3.12 of
          Ref.~\cite[lines 583--590]{Cirac2016MPDO_arXiv} states that every
          block of a canonical-form renormalization fixed point has a
          unit-modulus coefficient and an idempotent transfer map, and that
          the phase-class representatives have a joint residual isometry. The
          corresponding blueprint nodes are
          \path{thm:cpsv_weight_norm_one_and_block_rfp} and
          \path{thm:cpsv_phase_class_residual_isometries}.
\end{itemize}

Without the nonzero-coefficient convention, one could adjoin a normal tensor
with coefficient identically zero to any basis without changing any
positive-length matrix product vector. The cardinality conclusions of the
uniqueness statement, Theorem~2.10, and Lemma~A.5 of
Ref.~\cite{Cirac2016MPDO_arXiv} would then fail. The convention records the
source's intended reading and makes the formal definition consistent with
Equation~\texttt{eq:II\_CF1} of Ref.~\cite{Cirac2016MPDO_arXiv}, reproduced in
Eq.~\ref{eq:cpsvbnt_canonical_form_decomposition}.

\bibliographystyle{alpha}
\bibliography{references}

\end{document}
